\chapter{ML Metrics and HyperParameter Tuning}

\section{Performance Metrics}

Losses are for machines, while it is needed to define metrics for humans. Metrics are a measure of quality and can be exploited in order to choose the hyper-parameters of a model.

\subsection{Evaluation of supervised ML}

\begin{itemize}
    \item Efficiency: model building time, prediction time.
    \item Scalability: training set size, attribute number.
    \item Robustness: noise, missing data.
    \item Interpretability: model interpretability, model compactness.
\end{itemize}
These are independent from the correctness of the resulting $y'$ labels.

\subsection{Evaluation of Regression Techniques}

The quality of the prediction is assessed through loss functions such as mean squared error, mean absolute error and other ones.

\section{Evaluation of Classification Techniques}

The quality of the prediction can be assessed through:
\begin{itemize}
    \item Confusion Matrix.
    \item Accuracy.
    \item Precision, Recall, F-measure per class and on average.
    \item ROC curve and other metrics.
\end{itemize}

\subsection{Confusion Matrix}

The Confusion Matrix is used to evaluate a classification task implemented by binary classifiers. It is built as follows:
\begin{itemize}
    \item Rows: Actual Class
    \item Columns: Predicted Class
    \item TP: True Positives when $y = y'$ for the positive class.
    \item TN: True Negatives when $y = y'$ for the negative class.
    \item FP: False Positives when $y'$ is positive while $y$ is negative.
    \item FN: False Negatives when $y'$ is negative while $y$ is positive.
\end{itemize}

\begin{table}[h]
    \centering
    \begin{tabular}{c|c|c}
        \textbf{Actual} $\backslash$ \textbf{Predicted} & \textbf{P} & \textbf{N} \\ \hline
        \textbf{P} & True Positives (TP)  & False Negatives (FN) \\ \hline
        \textbf{N} & False Positives (FP) & True Negatives (TN) \\
    \end{tabular}
    \caption{Confusion Matrix}
\end{table}

It can go up to $(n \times n)$. The correctly predicted values are always on the diagonal.

\subsection{Accuracy}

Accuracy is the most widely-used metric for model evaluation and it is computed as follows:
\[
\text{Accuracy} = \frac{\text{Number of correctly classified objects}}{\text{Number of classified objects}}
\]
For binary classifiers:
\[
\text{Accuracy} = \frac{TP + TN}{TP + TN + FP + FN}
\]

In addition accuracy has the following properties:
\begin{itemize}
    \item Accuracy is related to the 0/1 loss: it goes as $1 -$ (0/1 loss value).
    \item Accuracy might be misleading in case of unbalanced classes or different class relevance and importance. Therefore, specific class measures are defined.
\end{itemize}

\subsection{Class Specific Measures}

Class specific measures need to be evaluated separately for each class $c$. They are the precision $p$, the recall $r$ and the F-measure, which is the harmonic mean of the precision and recall. They are defined as follows:
\begin{align*}
    \text{Precision}(c) \quad p &= \frac{\text{Number of objects correctly assigned to } c}{\text{Number of objects assigned to } c} \\[4pt]
    \text{Recall}(c) \quad r &= \frac{\text{Number of objects correctly assigned to } c}{\text{Number of objects belonging to } c} \\[4pt]
    \text{F-measure}(c) \quad F &= \frac{2rp}{r + p}
\end{align*}

For a binary classification problem the following are the simplified formulas for the \underline{positive} class $c$:
\begin{align*}
    \text{Precision} \quad p &= \frac{TP}{TP + FP} \\[4pt]
    \text{Recall (True Positive Rate)} \quad r &= \frac{TP}{TP + FN} \\[4pt]
    \text{F-measure} \quad F &= \frac{2rp}{r + p} = \frac{2TP}{2TP + FP + FN}
\end{align*}

\subsection{ROC Curve}

The Receiver Operating Characteristic (ROC) was developed for analyzing noisy signals and finding the trade-off between positive hits and false alarms. The ROC curve shows how the FP and TP rates change when changing the decision threshold. In $(0,0)$ we declare everything to be negative class, while in $(1,1)$ everything to be positive class. The diagonal line is random guessing if classes are balanced; below the diagonal line the prediction is opposite of the true class.\\
It is important to compute the area under the curve (AUC) and choose the model with the largest area.


M1 is better for small FPR, while M2 is better for large FPR.

\textit{How to build ROC curve?}
\begin{itemize}
    \item Use a classifier that produces the posterior probability $P(+|x)$ for each sample $x$.
    \item Apply a threshold at each unique value of $P(+|x)$.
    \item Compute TPR and FPR at each threshold.
\end{itemize}
Since in the example we have 10 samples, $m = 10$ columns. We need to order them according to probability horizontally.


\subsection{Learning Curve}
The Learning Curve shows how a metric, or loss, changes with varying training sample size. It is possible to use different losses for training and validation, and this enables the comparison of different ML methods. It is a tool to find out how much we benefit from adding more training data.

\section{HyperParameter Tuning}
Hyper-Parameters are parameters that are \textbf{not directly learnt within estimators}. \textbf{Hyper-Parameters define the model space} and they are passed as arguments to the constructor of the estimator classes in scikit-learn. The aim is to search the hyper-parameters space for the best cross-validation score metric. Hence, hyper-parameters are really important for the final model selection.\\
The \textbf{Validation Curve} can be used to assess the influence of a single hyper-parameter on the training score and the validation score. Hence, it is possible to find out whether the estimator is overfitting or underfitting for some hyper-parameter values.

\subsection{Grid Search}
Grid Search exhaustively generates candidates from a \textbf{grid of hyper-parameter values}. A different hypothesis is trained for all possible combinations of hyper-parameter values.

\section{Recap}
To sum up, these are the steps to follow when building an ML process:

\begin{itemize}
    \item Load/generate datasets.
    \item Choose samples, features and labels.
    \item Split the dataset in training, validation (or later cv) and testing.
    \item Pre-processing techniques.
    \item Define model and loss.
    \item Perform ERM with an algorithm (learn).
    \item Evaluate the obtained performance (metrics).
    \item Tune the hyper-parameters.
    \item Test the final hypothesis on test data.
    \item Use it in the wild for inference.
\end{itemize}
